\section{A jointly analytic parametrization of the glancing collar}
\label{joint:section}

The analyticity of the boundary does not by itself give analytic dependence
of its caustics on a transverse parameter. We prove that dependence by
constructing a conjugacy in which the transverse parameter is the angular
translation number. Rational periodicity supplies the compatibility
conditions needed to divide the two difference operators at their
resonances. Both the value and first-parameter-jet conditions use the
full collar of rational leaves; they do not follow from the
arbitrary-input action estimates alone.

By \Cref{thm:analytic,ana:analytic-action-interface}, we may use the
analytic normalized coordinate and action of
\Cref{ana:normalized-action}. Thus \(\theta=\phi(x)\),
\(r(\phi(x))\phi'(x)^3=c>0\), and \(D\) is real analytic near the real
diagonal, diagonally \(2\pi\)-periodic, and antisymmetric. On a fixed complex
neighborhood of that diagonal,
\begin{equation}
 D(x,x+b)=-\frac{b^3}{24}+O(b^5),\qquad
 D_1(x,x+b)=\frac{b^2}{8}+b^4 B(x,b),
 \label{joint:action-expansion}
\end{equation}
where \(B\) is holomorphic and periodic in \(x\). All domains below are
chosen inside this fixed neighborhood.

\begin{theorem}[Joint glancing parametrization]
\label{thm:joint-parametrization}
There are \(s,r_1>0\), with \(r_1<s\), and a real-symmetric holomorphic
function \(U(y,t)\) on
\[
 \{(y,t):|\Im y|<2s,\ |t|<r_1\}
\]
such that \(U-y\) is \(2\pi\)-periodic in \(y\), \(U\) is even in \(t\),
\(U(y,0)=y\), and
\begin{equation}
 D_1\bigl(U(y,t),U(y+t,t)\bigr)
 +D_2\bigl(U(y-t,t),U(y,t)\bigr)=0
 \label{joint:stationary-equation}
\end{equation}
for \(|\Im y|<s\), \(|t|<r_1\). Here real symmetry means
\(U(\bar y,\bar t)=\overline{U(y,t)}\).

Set
\begin{equation}
 \Theta(y,t)=\phi(U(y,t)),\qquad
 p(y,t)=-\mathcal S_1\bigl(\Theta(y,t),\Theta(y+t,t)\bigr).
 \label{joint:oriented-coordinates}
\end{equation}
On a possibly smaller product neighborhood, both functions are
holomorphic, real-symmetric, and even in \(t\). For positive small real
\(t\), the billiard map sends the line \((\Theta(y,t),p(y,t))\) to the
line \((\Theta(y+t,t),p(y+t,t))\). The variables \((y,t^2)\) give analytic
coordinates on a full one-sided phase collar. Moreover,
\begin{equation}
 dp\wedge d\Theta=A(t)\,dy\wedge dt,
 \qquad A(t)=\frac{c}{4}t+O(t^3),
 \label{eq:area-coefficient}
\end{equation}
where \(A\) is holomorphic and real on the real axis. In particular
\(t^2\) is a nondegenerate analytic first integral on that collar.
\end{theorem}

\subsection{Normalized equations and their linearization}

For a fixed parameter \(t\), write
\[
 T_zv(y)=v(y+z),\qquad
 \partial_t^{\mathrm{d}}v=\frac{T_{t/2}v-T_{-t/2}v}{t}.
\]
The superscript distinguishes this difference operator from an ordinary
parameter derivative. We abbreviate it to \(\partial\) when the parameter
is understood. A prime always denotes a \(y\)-derivative. For a candidate
\(U\), define
\begin{align}
 \mathcal R(U)&=D_1(U,T_tU)+D_2(T_{-t}U,U),\notag\\
 E(U)&=t^{-3}U'\mathcal R(U),\label{joint:normalized-residual}\\
 a(U)&=t^{-1}D_{12}(T_{-t/2}U,T_{t/2}U)
                    (T_{-t/2}U')(T_{t/2}U').
 \label{joint:coefficient-a}
\end{align}
The symbol \(a\) in this Section is the coefficient of the linearized
equation; the complex unit normal used in later Sections will be defined
separately.

\begin{lemma}[Regularity and automatic factorization]
\label{joint:factorization}
The expressions in \eqref{joint:normalized-residual}--\eqref{joint:coefficient-a}
extend holomorphically to \(t=0\), with
\begin{equation}
 E(U)(y,0)=\frac{(U')^2U''}{4},\qquad
 a(U)(y,0)=\frac{(U')^3}{4}.
 \label{joint:zero-parameter}
\end{equation}
If \(U\) is even in \(t\), so are \(E\) and \(a\). The mean of \(E\)
over a real period in \(y\) vanishes. For a variation \(\delta U=U'v\),
\begin{equation}
 E_U(U'v)=\partial(a\partial v)+(Ev)'.
 \label{joint:linearization}
\end{equation}
At an exact solution at a nonzero real parameter \(t_0\), assign an
arbitrary parameter jet \(\dot U=U'z\). Its total residual jet is
\begin{equation}
 \dot E=\partial\left(a\partial z+\frac{a}{t_0}\right).
 \label{joint:parameter-jet}
\end{equation}
All these identities hold wherever their shifted arguments are defined.
\end{lemma}

\begin{proof}
By antisymmetry,
\[
 \mathcal R=D_1(U,T_tU)-D_1(U,T_{-t}U).
\]
Put \(v_+=T_tU-U=tA_+\) and \(v_-=T_{-t}U-U=-tA_-\), where
\[
 A_\pm=\int_0^1 U'(y\pm\lambda t,t)\,d\lambda.
\]
The leading part of \(t^{-3}\mathcal R\) is
\[
 \frac18(A_++A_-)\frac{A_+-A_-}{t}.
\]
The last quotient is holomorphic, because
\[
 \frac{A_+-A_-}{t}
 =\int_0^1\int_{-\lambda}^{\lambda}U''(y+\mu t,t)\,d\mu\,d\lambda.
\]
The remaining part, using \eqref{joint:action-expansion}, is
\[
 t\left[A_+^4B(U,tA_+)-A_-^4B(U,-tA_-)\right].
\]
This proves regularity and the formula for \(E(y,0)\).
Differentiating \eqref{joint:action-expansion} in its second argument gives
\(D_{12}(x,x+b)=b/4+O(b^3)\), proving the formula for \(a\).
The same integral formulas show that, on smaller domains, the first two
variations of \(E\) are bounded by fixed polynomials in the appropriate
norms of the variations and their first two \(y\)-derivatives. There is no
negative power of \(t\) in these bounds.

Antisymmetry interchanges the forward and backward terms when \(t\) is
replaced by \(-t\), giving evenness. For real \(t\), the mean action
\[
 \frac{1}{2\pi}\int_0^{2\pi}D\bigl(U(y+\eta,t),U(y+\eta+t,t)\bigr)\,dy
\]
does not depend on \(\eta\). Differentiation at \(\eta=0\) gives
\(\langle U'\mathcal R\rangle=0\). Holomorphic continuation gives
\(\langle E\rangle=0\) for complex \(t\) as well.

For \eqref{joint:linearization}, subtract from the off-diagonal Hessian
terms those in which \(v\) is evaluated at \(y\). The remaining terms are
\[
 t^{-3}\left[c_+(T_tv-v)+c_-(T_{-t}v-v)\right],
\]
where
\[
 c_+=D_{12}(U,T_tU)U'T_tU',\qquad
 c_-=D_{12}(T_{-t}U,U)T_{-t}U'U'.
\]
Since \(T_{\pm t/2}a=c_\pm/t\), this is
\(\partial(a\partial v)\). The factored terms, including the variation
of the prefactor \(U'\), are \((Ev)'\).

Finally, at an exact solution the derivative of \(t^{-3}\) contributes
zero. Differentiating only the explicit shifts gives
\((c_+-c_-)/t_0^3=\partial a/t_0\). The assigned jet contributes
\(\partial(a\partial z)\) by \eqref{joint:linearization}. This proves
\eqref{joint:parameter-jet}.
\end{proof}

We specify the domains used to control complex shifts. For
\(\nu\in\{0,\tfrac12,1\}\), set
\begin{equation}
 \mathcal D^\nu(s,d)
 =\{(y,t):|t|<d,\ |\Im y|<s-\nu|\Im t|\}.
 \label{joint:slanted-domains}
\end{equation}
Fix \(d<s/2\), and choose both sufficiently small relative to the fixed
domain of \(D\). Also take \(d\) small enough that every positive real
step of size at most \(3d\) in the normalized coordinate gives a line
in the filled phase collar of \Cref{lem:rational-periodicity}. This is
possible uniformly around the boundary by the support-point formulas
and compactness. We work with candidates for which \(U-y\) and its first
three \(y\)-derivatives are uniformly small on \(\mathcal D^0(s,d)\).
The smallness constant is fixed once and for all. It ensures that every
composition above is defined, \(U'\) is close to one, and \(4a\) is close
to one on \(\mathcal D^{1/2}(s,d)\). Indeed, the leading term of \(a\)
is the product of three averaged first derivatives divided by four;
its remaining terms are \(O(d^2)\).

For \(0<\delta<d/100\), abbreviate
\[
 \mathcal D_j^\nu=\mathcal D^\nu(s-j\delta,d-j\delta),
 \qquad Y_j=\{y:|\Im y|<s-j\delta\}.
\]
All norms are supremum norms unless derivatives are indicated. Cauchy's
estimate between successive such domains costs a fixed power of
\(\delta^{-1}\). For example, polydisks of radius \(\delta/4\) in each
variable based in \(\mathcal D_{j+1}^\nu\) lie in
\(\mathcal D_j^\nu\). A shift by \(t/2\) increases the required domain
index \(\nu\) by \(1/2\).

\begin{proposition}[A Newton step across all retained resonances]
\label{prop:joint-newton}
There are constants \(C,B_0\), depending only on the fixed action domain
and the fixed smallness bounds above, with the following property.
Suppose \(U\) is real-symmetric, even in \(t\), and \(U-y\) is periodic,
and
\[
 \|E(U)\|_{\mathcal D^1_0}\le\eps<1.
\]
Let \(N\ge1\) satisfy \(e^{-N\delta}\le\eps\), and put
\(Q=1+N+\delta^{-1}\). If \(CQ^{B_0}\eps\) is sufficiently small, there
is a real-symmetric, even correction \(\Delta U\), periodic in \(y\),
such that
\begin{align}
 \max_{0\le j\le3}\|\partial_y^j\Delta U\|_{\mathcal D^0_{20}}
     &\le CQ^{B_0}\eps,\label{joint:step-correction}\\
 \|E(U+\Delta U)\|_{\mathcal D^1_{20}}
     &\le CQ^{B_0}\eps^2.
 \label{joint:step-error}
\end{align}
The exponent \(B_0\) is fixed; it is independent of \(N\), the retained
rationals, and the order of any formal approximation used to initialize
the iteration.
\end{proposition}

We prove the Proposition in three parts. In the estimates that follow,
\(P\) denotes a quantity bounded by \(CQ^B\), with a fixed exponent
\(B\). The exponent may increase a fixed finite number of times. At the
end we choose \(B_0\) larger than all these exponents, including those
needed for the stated smallness conditions.

\subsection{Rational corrections and compatible parameter jets}

\begin{lemma}[The two rational compatibility conditions]
\label{joint:rational-jets}
Under the hypotheses of \Cref{prop:joint-newton}, let
\(t_0=2\pi l/q\ne0\) be reduced, with \(q\le N\) and
\(|t_0|<d-2\delta\). There are holomorphic jets \(w,\dot w\) on
\(Y_{10}\), real-symmetric in \(y\), such that
\begin{align}
 E+\partial(a\partial w)&=O(P\eps^2),\notag\\
 \bigl(E+\partial(a\partial w)\bigr)^{\displaystyle\cdot}
       &=O(P\eps^2),\label{joint:rational-compatibility}\\
 w,\dot w&=O(P\eps).\notag
\end{align}
The dot differentiates the entire expression with respect to \(t\) at
\(t_0\), using the displayed assigned jet for \(w\). The bounds also
hold after any of the fixed derivative losses used below.
\end{lemma}

\begin{proof}
We first take \(t_0>0\). Write \(U_j(y)=U(y+jt_0,t_0)\) for
\(0\le j\le q\), and look for stationary nodes
\[
 \widetilde X_j=U_j+U'_j\xi_j,
 \qquad \xi_0=\xi_q=0.
\]
The endpoints differ by \(qt_0=2\pi l\). Scale the interior stationarity
equations by \(t_0^{-3}U'_j\). Their residual at \(\xi=0\) is
\(T_{jt_0}E\), and their linearization is a divergence matrix plus a
diagonal potential:
\begin{align}
 (A_0\xi)_j
 &=\frac{a_{j+1/2}(\xi_{j+1}-\xi_j)
              -a_{j-1/2}(\xi_j-\xi_{j-1})}{t_0^2},\notag\\
 V_j&=T_{jt_0}\left(E'-\frac{U''}{U'}E\right),\qquad
 a_{j+1/2}=a(y+(j+1/2)t_0,t_0).
 \label{joint:dirichlet-linearization}
\end{align}
All translations here are real and preserve the \(y\)-strip.

For completeness, the complex Dirichlet inverse has the same polynomial
bound as the real one. For \(A_0\xi=F\), introduce the fluxes
\(J_{j+1/2}=a_{j+1/2}(\xi_{j+1}-\xi_j)\). Then
\[
 J_{j+1/2}=J_{1/2}+t_0^2\sum_{i=1}^j F_i.
\]
The condition \(\sum_{j=0}^{q-1}(\xi_{j+1}-\xi_j)=0\) determines
\(J_{1/2}\). Since \(a_{j+1/2}^{-1}\) is uniformly close to 4,
\[
 \left|\sum_{j=0}^{q-1}a_{j+1/2}^{-1}\right|\ge c_0q.
\]
Consequently
\begin{equation}
 \|A_0^{-1}\|_{\infty\to\infty}\le Cq^2t_0^2\le CQ^2.
 \label{joint:dirichlet-inverse}
\end{equation}
Cauchy's estimate gives \(\|V\|_{Y_1}\le CQ\eps\). A Neumann
perturbation therefore gives the same inverse bound for \(A_0+V\) if
\(CQ^3\eps\) is small.

There is also a uniform polynomial bound for the nonlinear remainder.
The scaled equations contain only three adjacent nodes, fixed bounded
derivatives of \(D\), and the factor \(t_0^{-3}\). Since
\begin{equation}
 |t_0|^{-1}\le\frac{N}{2\pi},
 \label{joint:rational-lower-bound}
\end{equation}
their second derivatives in \(\xi\) have norm at most \(CQ^3\).
The fixed-point equation obtained with \((A_0+V)^{-1}\) thus preserves a
ball of radius \(CQ^2\eps\) and contracts there if \(CQ^7\eps\) is
small. This gives holomorphic corrections on \(Y_1\), with
\(\|\xi\|\le CQ^2\eps\). Uniqueness preserves periodicity and real
symmetry. Further derivatives follow by Cauchy estimates on subsequent
strips.

On real \(y\), the node increments are positive and small: their
uncorrected increments are bounded below by a fixed multiple of
\(t_0\), and the correction is smaller than that by
\eqref{joint:rational-lower-bound} and the polynomial smallness condition.
The stationary nodes therefore describe actual billiard trajectories
whose initial lines lie in the given phase collar. Such a line belongs
to one of the invariant caustic graphs. At the q-th step its angle has
increased by \(2\pi l\); invariance of the graph forces its momentum
also to return. Its rotation is \(l/q\). By
\Cref{lem:rational-periodicity}, this is the unique graph of that rotation,
and every point of it is q-periodic.

Use the local inverse of \(U(\,\cdot\,,t_0)\) near each \(U_j\) to write
\[
 \widetilde X_j(y)=U(\widetilde y_j(y),t_0),\qquad
 \widetilde y_j=y+jt_0+O(P\eps).
\]
These are holomorphic on \(Y_2\). A contraction for the displacement
from \(y+jt_0\), using \(U'\) close to one, gives this inverse without
any global inverse-domain assumption. On real points, the graph just
identified implies
\(\widetilde y_j=\widetilde y_1^{\,j}\) as iterates and
\(\widetilde y_q=y+qt_0\). Hence
\[
 G(y)=\frac1q\sum_{j=0}^{q-1}(\widetilde y_j(y)-jt_0)
\]
satisfies \(G(\widetilde y_1(y))=G(y)+t_0\). It is close to the identity,
and Cauchy's estimate makes \(G'\) close to one on \(Y_3\). Its inverse
exists on \(Y_4\), again by contraction. The function
\begin{equation}
 U_*(y)=U(G^{-1}(y),t_0)
 \label{joint:exact-rational-parametrization}
\end{equation}
is an exact solution at \(t_0\), and
\(U_*-U(\,\cdot\,,t_0)=O(P\eps)\). The identities first obtained on
the real axis extend throughout the stated complex strips.

We next assign a compatible parameter jet to \(U_*\). Start with
\(\dot U_*=U_t(\,\cdot\,,t_0)\). Its total residual jet is
\(O(P\eps)\) on \(Y_6\). Indeed \(E_t=O(P\eps)\) by Cauchy's
estimate, and replacing \(U\) by \(U_*\), with the same assigned jet,
changes the explicit differentiated formula for \(E\) by
\(O(P\eps)\). Only finitely many derivatives occur, and every
denominator \(t_0\) is controlled by \eqref{joint:rational-lower-bound}.
More explicitly, for any assigned jet \(V\), the raw derivative is
\begin{align}
 \dot E={}&-3t^{-4}U'\mathcal R+t^{-3}V'\mathcal R\notag\\
 &+t^{-3}U'\bigl[
 D_{11}(U,T_tU)V
 +D_{12}(U,T_tU)(T_tV+T_tU')\notag\\
 &\hspace{35mm}
 +D_{21}(T_{-t}U,U)(T_{-t}V-T_{-t}U')
 +D_{22}(T_{-t}U,U)V\bigr].
 \label{joint:raw-parameter-derivative}
\end{align}
This formula makes both the fixed derivative count and the worst
explicit denominator \(t^{-4}\) transparent.

Let
\[
 \Pi v=\frac1q\sum_{j=0}^{q-1}T_{jt_0}v
\]
be the orbit-average projection. Equation \eqref{joint:parameter-jet}
shows that \(\Pi\dot E_*=0\). This uses only
\(\partial=T_{-t_0/2}(T_{t_0}-1)/t_0\) and \(\Pi T_{t_0}=\Pi\);
the half-shift need not act trivially. To write the cyclic inverse
explicitly, put \(T=T_{t_0}\) and
\[
 R_q=\frac1q\sum_{j=0}^{q-1}jT^j.
\]
Since \(T^q=1\), one has \((T-1)R_q=1-\Pi\). Thus on \(\ker\Pi\),
\(\partial^{-1}=t_0R_qT_{t_0/2}\), and
\[
 \|\partial^{-1}\|\le Cq|t_0|.
\]
The polynomial smallness condition makes \(U_*\) close enough to \(U\)
through the required derivatives that \(4a_*\) is still close to one.
The operator \((1-\Pi)a_*\) on \(\ker\Pi\) is therefore invertible: its
difference from multiplication by \(1/4\) has norm at most
\(2\|a_*-1/4\|\). Factoring through these
three operators proves
\begin{equation}
 \left\|(\partial a_*\partial)^{-1}\right\|_{\ker\Pi}
       \le Cq^2t_0^2.
 \label{joint:rational-jet-inverse}
\end{equation}
Correcting \(\dot U_*\) by \(U_*'z\), with \(z\in\ker\Pi\), now makes
\(\dot E_*=0\), at cost \(O(P\eps)\).

Define \(w=(U_*-U)/U'\) and assign its derivative by differentiating
this quotient with the corrected jet of \(U_*\). Both \(w\) and
\(\dot w\) are \(O(P\eps)\). Taylor expansion and
\eqref{joint:linearization} give
\[
 0=E(U_*)=E+\partial(a\partial w)+(Ew)'+O(P\eps^2).
\]
Differentiating this Taylor identity once in the assigned parameter jets
gives the second assertion in \eqref{joint:rational-compatibility}.
The term \((Ew)'\) and its jet are quadratic, because \(E,E_t,w,\dot w\)
are all bounded by a polynomial times \(\eps\). The same is true of
the differentiated Taylor remainder.

Here is a concrete accounting for the fixed-polynomial assertion. The
mesh inverse and correction cost at most \(CQ^2\); the nonlinear mesh
Hessian costs \(CQ^3\). Using subsequent strip losses for the fixed
number of derivatives, comparison of the initially assigned residual
jet costs at most \(CQ^{10}\eps\), and its correction costs at most
\(CQ^{12}\eps\). In a first parameter derivative of a quadratic
remainder, either one correction is replaced by its assigned jet, or
one coefficient is differentiated. The latter introduces at most one
additional \(t_0^{-1}\), fixed derivatives of \(D\), and bounded
\(U_t,U'_t\). Before further y-derivatives, the value remainder is
bounded by \(CQ^8\eps^2\) and its first jet by \(CQ^{18}\eps^2\): a
correction and its first y-derivative cost \(Q^2\eps,Q^3\eps\), whereas
the assigned correction jet and its first y-derivative cost
\(Q^{12}\eps,Q^{13}\eps\). Multiplication by \(t_0^{-3}\) gives these
bounds; differentiating that factor or a coefficient costs less.
Allowing six additional y-derivatives on a fixed smaller strip therefore
bounds both remainders by \(CQ^{24}\eps^2\), more than is needed below.
Enlarging
to \(CQ^{40}\) covers all bounds and smallness conditions in this
Lemma. These exponents are deliberately inessential upper bounds; in
particular they do not depend on any formal approximation order.

For negative \(t_0\), apply the positive construction at \(-t_0\) and
use the even symmetry, reversing the sign of the assigned parameter jet.
\end{proof}

\subsection{Division at the resonant parameters}

For a periodic function \(v\), write
\[
 v(y,t)=\sum_{n\in\Z}v_n(t)e^{\ii ny},\qquad
 d_n(t)=\frac{2\ii\sin(nt/2)}{t}.
\]
Thus \((\partial v)_n=d_nv_n\), with \(d_n(0)=\ii n\).
The nonzero zeros of \(d_n\) are simple and occur at
\(t_j=2\pi j/n\), \(j\ne0\).

\begin{lemma}[Buffered Fourier division]
\label{joint:fourier-division}
Let \(\nu\in\{1/2,1\}\), and suppose a mean-zero periodic function
\(F\) is holomorphic on \(\mathcal D_{j_0}^\nu\), with norm at most
\(M\). Suppose, for \(1\le |n|\le N\), its values at the divisor zeros
\(0<|t_j|<d-b\delta\), where \(b\ge j_0+1\), satisfy
\[
 |F_n(t_j)|\le\eta e^{-|n|(s-a\delta)},
 \qquad a\ge j_0.
\]
Define, for these retained modes,
\begin{align}
 R_n(t)&=\sum_{0<|t_j|<d-b\delta}F_n(t_j)
       \frac{\sin(n(t-t_j)/2)}{n(t-t_j)/2},\label{joint:sinc-interpolant}\\
 H_n(t)&=\frac{F_n(t)-R_n(t)}{d_n(t)},
 \label{joint:divided-coefficient}
\end{align}
and put \(H_0=0\), \(H_n=0\) for \(|n|>N\).
On \(|t|<d-(b+1)\delta\), every retained \(H_n\) is holomorphic.
For an integer \(J\) larger than \(a\), \(j_0\), and \(b+1\),
\begin{align}
 \|H\|_{\mathcal D_J^{\nu-1/2}}&\le P(M+\eta),\notag\\
 \|\partial H-F\|_{\mathcal D_J^\nu}
       &\le P\bigl(\eta+Me^{-N\delta}\bigr).
 \label{joint:division-bounds}
\end{align}
Fixed derivatives are bounded after fixed additional domain losses by
increasing \(P\). Real symmetry and evenness in \(t\) are preserved.
\end{lemma}

\begin{proof}
At two distinct divisor zeros, the sinc kernel in
\eqref{joint:sinc-interpolant} is zero; at its own center it is one.
Thus \(R_n(t_j)=F_n(t_j)\), so the zeros in the quotient are removable.
There are at most \(CN\) interpolation points for each mode, because
\(d<1\). The identity
\[
 \frac{\sin z}{z}=\int_0^1\cos(uz)\,du
\]
gives
\begin{equation}
 |R_n(t)|+N^{-1}|R_n'(t)|
 \le P\eta e^{-|n|(s-a\delta-|\Im t|/2)}.
 \label{joint:interpolant-bound}
\end{equation}
Outside disks centered at the divisor zeros with radius
\(c_1\min(|n|^{-1},\delta)\), elementary estimates for the sine give
\[
 |d_n(t)|^{-1}\le Pe^{-|n\Im t|/2}.
\]
Here and below a harmless fixed constant can be absorbed into \(P\).
For \(|nt|\le1\), the function \(\sin(nt/2)/(nt/2)\) is bounded
away from zero. Thus \(|d_n(t)|^{-1}\le C/|n|\), which has the
required bound since \(e^{-|n\Im t|/2}\ge e^{-1/2}\).
On the remaining region, write \(nt/2=u+\ii v\) and use
\(|\sin(u+\ii v)|^2=\sin^2u+\sinh^2v\). Outside the nonzero-root
disks, \(|\sin(u+\ii v)|\) is bounded below by a fixed multiple of
\(\min(1,|n|\delta)e^{|v|}\): if the nearest sine zero is zero,
the condition \(|nt|>1\) supplies the required fixed separation.
Multiplication by \(|t|<1\) then gives a polynomial bound in \(Q\).
Inside the nonzero-root disks, apply maximum modulus
to the holomorphic quotient. The change in any exponential weight is
bounded by \(e^{C|n|\min(|n|^{-1},\delta)}\), a fixed constant.
The parameter buffer ensures that every disk meeting the target domain
is contained in the domain of \(F_n\) and has its center among the
interpolation points.

The input strip gives
\[
 |F_n(t)|\le M e^{-|n|(s-j_0\delta-\nu|\Im t|)}.
\]
After division its imaginary-parameter loss is
\((\nu-1/2)|\Im t|\). The interpolant contribution after division has
no such loss. Hence
\[
 |H_n(t)|\le P(M+\eta)
 e^{-|n|(s-a\delta-(\nu-1/2)|\Im t|)}.
\]
Summing on the smaller y-strip proves the first bound, with a geometric
sum of size at most \(C/\delta\). The retained-mode residual is
\(-R_n\), and the omitted tail of \(F\) gains at least
\(e^{-N\delta}\); summing similarly proves the second bound. Each
fixed y-derivative costs only a fixed power of \(N\) or \(\delta^{-1}\).
Parameter derivatives follow from the explicit bound on \(R_n'\) or
from Cauchy's estimate on a further buffered domain.

For real-symmetric input, conjugate modes have conjugate output, because
\(d_{-n}(\bar t)=\overline{d_n(t)}\). If the input is even, the divisor
zeros occur in opposite pairs and the interpolant is even; so is the
quotient. These properties survive the symmetric truncation.
\end{proof}

\begin{proof}[Proof of \Cref{prop:joint-newton}]
First divide \(f=-E\). Its mean is zero by
\Cref{joint:factorization}. For \(|n|\le N\), every nonzero zero of
\(d_n\) has reduced denominator at most \(N\). At those zeros with
\(|t_j|<d-3\delta\), the first equation of
\eqref{joint:rational-compatibility} gives
\[
 |f_n(t_j)|\le P\eps^2e^{-|n|(s-10\delta)}.
\]
Apply \Cref{joint:fourier-division} with
\(\nu=1\), \(j_0=0\), \(a=10\), \(b=3\), and \(J=11\).
Call its output \(h\). Then
\begin{equation}
 h=O(P\eps)\ \hbox{on }\mathcal D_{11}^{1/2},\qquad
 \partial h+E=O(P\eps^2)\ \hbox{on }\mathcal D_{11}^{1}.
 \label{joint:first-division}
\end{equation}
The term involving the interpolation values is absorbed by the
smallness condition. Define
\begin{equation}
 b=a^{-1}(h+\lambda),\qquad
 \lambda(t)=-\frac{\langle a^{-1}h\rangle}{\langle a^{-1}\rangle}.
 \label{joint:mean-adjustment}
\end{equation}
The denominator is bounded away from zero because \(a^{-1}\) is close
to four. Thus \(b\) has mean zero and is \(O(P\eps)\) on
\(\mathcal D_{12}^{1/2}\).

We need new compatibility conditions before dividing \(b\). Fix a
retained rational \(t_0\) with \(|t_0|<d-14\delta\), and put
\(g=a\partial w\), using \Cref{joint:rational-jets}. For retained
nonzero modes the exact first-division identity is
\[
 d_n(h_n-g_n)=-(E+\partial g)_n-R_n.
\]
Its value and its first parameter derivative at \(t_0\) are bounded
by
\(P\eps^2e^{-|n|(s-10\delta)}\), by the two compatibility conditions
and \eqref{joint:interpolant-bound}. If \(d_n(t_0)\ne0\), reducedness
of \(l/q\) gives
\[
 |\sin(\pi nl/q)|\ge c/q,
 \qquad |d_n(t_0)|^{-1}\le Cq|t_0|.
\]
If \(d_n(t_0)=0\), then
\[
 |d_n'(t_0)|=|n|/|t_0|,
\]
and differentiation of the product gives the same bound for
\(h_n-g_n\). In particular, this comparison does not require the
resonant modes of \(g\) to vanish.

For the omitted modes, \(h_n=0\), whereas the analytic function \(g\)
has norm \(P\eps\) on \(Y_{10}\). Its tail on \(Y_{13}\) is bounded
by \(P\eps e^{-3N\delta}\), hence by \(P\eps^2\). Summing retained
and omitted modes gives
\[
 h-(g-g_0)=O(P\eps^2)\quad\hbox{on }Y_{13}.
\]
Since \(\langle a^{-1}g\rangle=\langle\partial w\rangle=0\), the mean
choice in \eqref{joint:mean-adjustment} gives
\(\lambda=g_0+O(P\eps^2)\). Consequently
\begin{equation}
 b-\partial w=O(P\eps^2)\quad\hbox{on }Y_{13}.
 \label{joint:second-compatibility}
\end{equation}
At a nonzero divisor zero for a retained mode, \((\partial w)_n=0\).
Thus \eqref{joint:second-compatibility} supplies the values needed for
the second division, with bound
\(P\eps^2e^{-|n|(s-13\delta)}\).

Apply \Cref{joint:fourier-division} to \(b\), now with
\(\nu=1/2\), \(j_0=12\), \(a=13\), \(b=15\), and \(J=17\).
To avoid ambiguity, the letter \(b=15\) here denotes the integer buffer
in that Lemma, not the function in \eqref{joint:mean-adjustment}.
Call the output \(v\). The remaining half of the imaginary-parameter
loss is canceled, giving
\begin{equation}
 v=O(P\eps)\ \hbox{on }\mathcal D_{17}^{0},\qquad
 \partial v-b=O(P\eps^2)\ \hbox{on }\mathcal D_{17}^{1/2}.
 \label{joint:second-division}
\end{equation}
Both divisions, multiplication by \(a^{-1}\), and the mean adjustment
preserve real symmetry and evenness.

Because \(ab=h+\lambda(t)\) and \(\partial\lambda=0\),
\eqref{joint:first-division} and \eqref{joint:second-division} imply
\[
 E+\partial(a\partial v)=O(P\eps^2)
       \quad\hbox{on }\mathcal D_{18}^{1}.
\]
Here application of \(\partial\) is bounded after the indicated strip
loss: use
\(\partial F=\int_{-1/2}^{1/2}F'(y+\mu t,t)\,d\mu\).
Take \(\Delta U=U'v\). Cauchy's estimate on the remaining domains gives
the required bounds through three y-derivatives. Equation
\eqref{joint:linearization} gives a quadratic linearized residual,
since \((Ev)'=O(P\eps^2)\). The normalized integral formulas in the
proof of \Cref{joint:factorization} give
\[
 E(U+\Delta U)-E(U)-E_U(\Delta U)=O(P\eps^2)
       \quad\hbox{on }\mathcal D_{20}^{1},
\]
uniformly at \(t=0\) as well.

Every cost used in this proof is a fixed number of additions,
multiplications, bounded inverses, Cauchy losses, sums over at most
\(CN\) roots per mode, and Fourier sums bounded by a power of
\(\delta^{-1}\). The rational part has the fixed ledger given in
\Cref{joint:rational-jets}. The two divisions add only finitely many
powers of \(Q\), and the normalized remainder adds a fixed derivative
cost. Choosing one larger exponent \(B_0\) proves
\eqref{joint:step-correction}--\eqref{joint:step-error}, with all required
smallness conditions included. No step involves a number of derivatives
depending on \(N\) or on the initialization order.
\end{proof}

\subsection{Initialization and convergence}

\begin{lemma}[Arbitrarily accurate formal initial data]
\label{joint:formal-initialization}
For every fixed integer \(m\ge1\), there are real analytic periodic
functions \(u_1,\ldots,u_m\) such that
\[
 U^{(m)}(y,t)=y+\sum_{j=1}^m u_j(y)t^{2j}
\]
satisfies \(E(U^{(m)})=O(t^{2m+2})\). All coefficients are holomorphic
on a common fixed smaller y-strip. For \(d_0\) sufficiently small,
\(U^{(m)}\) satisfies the fixed smallness conditions with slack and
\(\|E(U^{(m)})\|\le C_m d_0^{2m+2}\) on its initial shifted domain.
\end{lemma}

\begin{proof}
The coefficient of \(t^{2j}\) involving the newly chosen \(u_j\) is
\(u_j''/4\), by \eqref{joint:zero-parameter}; all other terms at that
order depend only on earlier coefficients. Their mean is zero, because
the full residual has mean zero for every candidate. Solve the periodic
equation for \(u_j\), choosing its mean to be zero. Integration, or
division of its nonzero Fourier modes by \(-n^2\), preserves
holomorphy on every strictly smaller strip and real symmetry. At every
fixed order the forcing uses only finitely many derivatives of
holomorphic functions on the original open strip. We may therefore fix
one smaller strip first and carry out any finite number of recursions
there, with constants allowed to depend on \(m\).

Evenness is built into the construction. On the fixed smaller strip,
the first three y-derivatives of \(U^{(m)}-y\) are \(O(C_m d_0^2)\),
and the residual bound follows from its first omitted power. Taking
\(d_0\) smaller gives all required slack.
\end{proof}

\begin{proof}[Construction of the function in \Cref{thm:joint-parametrization}]
Choose a fixed sufficiently small y-width \(s_0\). The constants in
\Cref{prop:joint-newton} are now fixed. Start with
\Cref{joint:formal-initialization}, choosing its order below, and let
the initial parameter radius be \(d_0<s_0/4\). At step \(i\), use losses
\[
 \delta_i=2^{-i}d_0/200,
\]
and reduce both current radii by \(20\delta_i\). The limiting parameter
radius is at least \(4d_0/5\), and the limiting y-width is at least
\(s_0-d_0/5\). If the current error is \(\eps_i>0\), choose
\[
 N_i=\left\lceil\delta_i^{-1}\log(1/\eps_i)\right\rceil.
\]
If it is zero, the construction has already finished. Enlarge one fixed
integer \(B\) if necessary. Since every power of
\(1+\log(1/\eps_i)\) is bounded by a constant times \(\eps_i^{-1/4}\),
the step estimates imply
\begin{align}
 \eps_{i+1}&\le C d_0^{-B}2^{Bi}\eps_i^{7/4},
 \label{joint:error-iteration}\\
 \|\Delta U_i\|_{C_y^3}
     &\le C d_0^{-B}2^{Bi}\eps_i^{3/4}.
 \label{joint:correction-iteration}
\end{align}
The norms in the second line are on the next domains.

Choose first \(m\) so that \(2m+2>4B+4\), then \(d_0\) sufficiently
small. Its initial error is at most \(d_0^{4B+4}\). Direct substitution
in \eqref{joint:error-iteration} shows inductively that
\begin{equation}
 \eps_i\le d_0^{4B+4}2^{-(4B+4)i}.
 \label{joint:error-majorant}
\end{equation}
Indeed the ratio of the resulting right-hand side of
\eqref{joint:error-iteration} to the next bound in
\eqref{joint:error-majorant} is at most
\(C d_0^{2B+3}2^{4B+4-(2B+3)i}\), which is less than one for small
\(d_0\). Equation \eqref{joint:correction-iteration} then bounds the
corrections by
\[
 C d_0^{2B+3}2^{-(2B+3)i}.
\]
They are summable and preserve the fixed smallness hypotheses with
slack. The same estimates ensure all polynomial smallness conditions
needed at every step. This order of choices is important: the exponent
\(B\) is fixed before \(m\), whereas the possibly large constant
\(C_m\) is absorbed by choosing \(d_0\) afterward.

The corrections converge uniformly, together with their first three
y-derivatives, on the limiting domains. Their limit is holomorphic,
real-symmetric, even, and periodic modulo \(y\). The residual tends to
zero on the limiting \(\mathcal D^1\), so the limit solves
\eqref{joint:stationary-equation}. At \(t=0\),
\eqref{joint:zero-parameter} and \(U'\ne0\) give \(U''=0\).
Periodicity forces \(U(y,0)=y+c_0\), with a real constant \(c_0\).
Replacing \(y\) by \(y-c_0\) removes this constant without changing
any of the other properties. Finally choose a smaller product domain
with enough y-margin for the two shifts. This gives the first assertion
of \Cref{thm:joint-parametrization}.
\end{proof}

\subsection{The phase collar and its area form}

\begin{proof}[Completion of the proof of \Cref{thm:joint-parametrization}]
The change from \(\mathcal S\) to \(D\) subtracted a telescoping
endpoint integral and changed the boundary coordinate. Thus
\eqref{joint:stationary-equation} is equivalent to
\[
 \mathcal S_1(\Theta,T_t\Theta)
       +\mathcal S_2(T_{-t}\Theta,\Theta)=0.
\]
The momentum formula in \eqref{joint:oriented-coordinates} therefore
matches at successive vertices, giving the claimed billiard dynamics.
All steps are positive and small for positive small real \(t\).

The angle \(\Theta\) is even. Antisymmetry of \(\mathcal S\) and the
last stationarity equation give
\[
 p(y,-t)=-\mathcal S_1(\Theta,T_{-t}\Theta)
       =\mathcal S_2(T_{-t}\Theta,\Theta)=p(y,t).
\]
Hence both oriented line coordinates are jointly analytic and even.
Taylor expansion of the support-point formula for the momentum gives
\begin{equation}
 p(y,t)=h(\Theta(y,t))
        -\frac18 r(\Theta(y,t))\phi'(y)^2t^2+O(t^4).
 \label{joint:collar-expansion}
\end{equation}
For example, if \(\beta=T_t\Theta-\Theta\), that formula first gives
\(p=h(\Theta)-r(\Theta)\beta^2/8+O(\beta^3)\). Since
\(\beta=\phi'(y)t+O(t^2)\), it determines the quadratic term;
evenness then removes the cubic term from the full expression.

An even holomorphic function of \(t\) is holomorphic in \(t^2\).
The derivative of \(\Theta\) with respect to \(y\) at zero is
\(\phi'(y)>0\), and the coefficient of \(t^2\) in the inward momentum
displacement in \eqref{joint:collar-expansion} is strictly positive.
The analytic inverse function theorem, uniformly around the compact
real period circle, consequently gives local collar coordinates.
They are globally one-to-one modulo the period after shrinking the
collar: for each fixed small real \(t^2\), \(\Theta_y>0\) determines a
unique \(y\) from the angle, and at a fixed angle the inward momentum
displacement has strictly positive derivative in \(t^2\).
Thus the coordinates give a full one-sided phase collar, and the shift preserves the
nondegenerate analytic coordinate \(t^2\).

Write the pulled-back area form as
\(B_0(y,t)\,dy\wedge dt\). It is holomorphic and periodic in \(y\).
The billiard map preserves \(dp\wedge d\Theta\), and in these
coordinates it is \((y,t)\mapsto(y+t,t)\). Therefore
\(B_0(y+t,t)=B_0(y,t)\). For every small real \(t\) whose ratio to
\(2\pi\) is irrational, density of the translation orbit makes
\(B_0(\,\cdot\,,t)\) constant. Its nonzero Fourier coefficients
vanish on this set of parameters and hence identically by holomorphy.
Thus \(B_0(y,t)=A(t)\).

Finally, \eqref{joint:collar-expansion} gives the coefficient of
\(dy\wedge dt\) as
\[
 \frac14 r(\phi(y))\phi'(y)^3t+O(t^3)
       =\frac{c}{4}t+O(t^3).
\]
The coefficient is odd because the line coordinates are even. This
proves \eqref{eq:area-coefficient} and completes the theorem.
\end{proof}
